%=============================================================
%  Auditing Label- and Group-Conditional Coverage for Student Dropout Prediction
%  A Class-Conditional Conformal Prediction Approach
%
%  Rimsha Bibi, M. Umair Shahzad
%  Department of Mathematics, University of Okara, Pakistan
%
%  ------------------------------------------------------------
%  HOW TO COMPILE
%  ------------------------------------------------------------
%  Upload this .tex file and the two Figure4 PNG files together.
%  The existing figures are drawn using pgfplots; the new
%  Figure 4A and Figure 4B are supplied as separate PNG files.
%
%  Overleaf: New Project > Blank Project, paste this whole file
%  into main.tex, and Recompile. Engine: pdfLaTeX.
%  The first compile is slower because the plots are drawn.
%
%  ------------------------------------------------------------
%  JOURNAL FORMATTING
%  ------------------------------------------------------------
%  This file uses the plain `article' class on purpose, so it
%  compiles under a plain article template.
%  Format for a journal only AFTER the journal is chosen:
%  copy everything between \begin{document} and \end{document}
%  into that publisher's template. Nothing in the body below
%  depends on the article class.
%
%  ------------------------------------------------------------
%  BEFORE SUBMISSION
%  ------------------------------------------------------------
%  Confirm every entry in the bibliography against the
%  publisher's own page: volume, issue, year, page range.
%=============================================================
\documentclass[11pt,a4paper]{article}

\usepackage[utf8]{inputenc}
\usepackage[T1]{fontenc}
\usepackage[margin=2.5cm]{geometry}
\usepackage{amsmath,amssymb,amsthm}
\usepackage{booktabs}
\usepackage{graphicx}
\usepackage{caption}
\usepackage{subcaption}
\usepackage{siunitx}
\usepackage{tikz}
\usepackage{pgfplots}
\pgfplotsset{compat=1.18}

% figure palette -- checked for colour-blind separation
\definecolor{fblue}{HTML}{2A78D6}
\definecolor{forange}{HTML}{EB6834}
\definecolor{faqua}{HTML}{1BAF7A}
\definecolor{fgrey}{HTML}{8C8C86}
\definecolor{fdim}{HTML}{9AA4AE}

\usepackage[colorlinks=true,linkcolor=black,citecolor=black,urlcolor=blue]{hyperref}

% --- float placement: keeps figures/tables from leaving large gaps ---
\renewcommand{\topfraction}{0.9}
\renewcommand{\bottomfraction}{0.8}
\renewcommand{\textfraction}{0.07}
\renewcommand{\floatpagefraction}{0.75}
\setcounter{topnumber}{3}
\setcounter{bottomnumber}{2}
\setcounter{totalnumber}{5}

\theoremstyle{plain}
\newtheorem{theorem}{Theorem}
\newtheorem{lemma}[theorem]{Lemma}
\theoremstyle{definition}
\newtheorem{definition}[theorem]{Definition}
\theoremstyle{remark}
\newtheorem{remark}[theorem]{Remark}

\newcommand{\Prob}{\mathbb{P}}
\newcommand{\Ind}{\mathbf{1}}
\newcommand{\Yset}{\mathcal{Y}}
\newcommand{\qhat}{\hat{q}}
\newcommand{\pihat}{\hat{\pi}}

\title{\bfseries Auditing Label- and Group-Conditional Coverage in Three-Class Student Outcome Prediction}

\author{%
Rimsha Bibi \\
\small Department of Mathematics, University of Okara, Okara, Pakistan \\
\small \texttt{rimsharao000@gmail.com}
\and
Dr. M. Umair Shahzad\thanks{Corresponding author: \texttt{mushahzad@uo.edu.pk}} \\
\small Department of Mathematics, University of Okara, Okara, Pakistan
}
% SUBMISSION CHECK: confirm applicable ethics policy and add repository URL and DOI after deposit. Check target journal declaration requirements.
\date{}

\begin{document}
\maketitle

%=============================================================
\begin{abstract}
We investigate whether near-nominal marginal conformal coverage masks
uneven uncertainty across student outcomes in the 4,424-record UCI student
attrition benchmark. A random-forest classifier was evaluated using 100
stratified training, calibration, and test partitions at 90\% nominal
coverage. Standard split conformal prediction achieved 90.25\% mean
marginal coverage but only 71.73\% for students labelled \textsf{Enrolled};
the latter fell below 90\% in every partition. Label-Mondrian conformal
prediction raised \textsf{Enrolled} coverage to 90.27\%, increasing mean
prediction-set size from 1.427 to 1.606 labels. A further audit of five
student attributes identified disparities within outcome--group cells:
for \textsf{Enrolled} scholarship holders, mean coverage increased from
60.79\% under split conformal prediction to 91.46\% under group-by-label
calibration. Small calibration cells sometimes produced large, less
informative sets. Results were examined under alternative classifiers,
feature subsets, calibration sizes, and score constructions. In a
100-partition matched automatic-decision comparison, Mondrian singleton
predictions had 90.72\% accuracy, below the 93.32\% obtained by selecting
an equal number of high-confidence predictions using maximum model
probability. These results indicate that educational prediction systems
should report marginal, outcome-conditional, and subgroup coverage
alongside prediction-set size and matched selection performance, rather
than rely on a single aggregate coverage statistic.
\end{abstract}

\noindent\textbf{Keywords:} conformal prediction; class-conditional coverage;
Mondrian conformal prediction; student dropout; educational data mining;
distribution-free uncertainty quantification; selective prediction

\vspace{1em}

%=============================================================
\section{Introduction}

Student dropout is among the most consequential and most studied outcomes in
higher education. It carries a direct cost to the student, a financial cost to
the institution, and a measurable social cost, and the prospect of identifying
at-risk students early enough for intervention has made it one of the standard
proving grounds of educational data mining. A large literature now reports
classifiers that predict, from administrative and first-year academic records,
whether a student will withdraw, remain enrolled, or graduate.

That literature is mature on the question of predictive accuracy. On the
Portalegre cohort in particular---4{,}424 students at a Portuguese polytechnic,
released publicly by Realinho et al.~\cite{realinho2022} and now a common
benchmark---recent work reports carefully validated results. Gu and
Zhang~\cite{gu2026} compare eight classifiers with SHAP-based interpretability
and report a macro-AUC of 0.891 with a 95\% confidence interval of
$[0.876, 0.902]$. Guamán-Hidalgo and Enciso~\cite{guaman2026} study student outcome
prediction using a three-class formulation; their reported performance is not
directly comparable without matched preprocessing, model selection, and splits. Elmasry et
al.~\cite{elmasry2026} go further and attach split conformal prediction to the
same data, reporting that the conformal layer ``achieves empirical coverage
close to the desired confidence level, ensuring reliable uncertainty
estimates''.

The studies quantify different things, and the distinction matters. A
confidence interval on an AUC, as in~\cite{gu2026}, is a statement about a
population parameter: it says how precisely the model's average discriminative
ability has been estimated across the cohort. It says nothing about any
individual student. Conformal prediction, as in~\cite{elmasry2026}, does supply
a per-student object---a set of labels containing the truth with a stated
probability---but the guarantee attached to that object is \emph{marginal}:
\begin{equation}
\Prob\big(Y \in C(X)\big) \;\ge\; 1-\alpha ,
\label{eq:marginal-intro}
\end{equation}
where the probability averages over the independent calibration sample and the future test pair $(X,Y)$ (conditional on the fitted predictor). What an advisor
implicitly assumes on being told that a prediction set is ``90\% reliable'' is
the stronger, conditional statement
\begin{equation}
\Prob\big(Y \in C(X) \,\big|\, Y = c\big) \;\ge\; 1-\alpha
\qquad \text{for each outcome } c ,
\label{eq:conditional-intro}
\end{equation}
and \eqref{eq:marginal-intro} does not imply \eqref{eq:conditional-intro}. The
common marginal threshold need not meet the nominal coverage target in every
outcome class, as the present empirical audit demonstrates.

This study measures, on the Portalegre cohort, the previously established
failure of marginal coverage to ensure class-conditional coverage
\cite{ding2023,lofstrom2015}. It is a benchmark-specific audit, not a new
conformal validity theorem.

To our knowledge, this is the first combined class- and group-by-label
conformal coverage audit on this specific student-outcome benchmark to
include an equally sized, matched maximum-probability selection baseline.
The contribution is the empirical reporting protocol rather than a new
validity result.

Our empirical contributions are:
\begin{enumerate}
\item A 100-partition class-wise coverage audit showing that standard split CP
      undercovers \textsf{Enrolled} in 100/100 partitions (mean 71.73\%).
\item A comparison with Mondrian CP, whose 100-partition class-wise mean
      coverages are 90.59\%, 90.27\% and 90.33\%, together with set-size costs.
\item A five-attribute group and group-by-label coverage audit, plus a
      transparently labelled two-cluster score-distribution adaptation.
\item Fixed-training calibration-size and within-training hyperparameter
      tuning studies, alongside score, model and feature-availability checks.
\item A matched-auto-decision comparison showing that ordinary maximum
      probability outperforms Mondrian singleton decisions on this dataset;
      no conditional risk guarantee is claimed for singletons.
\end{enumerate}


The remainder of the paper is organised as follows.
Section~\ref{sec:related} places the work against the dropout-prediction and
conformal-prediction literatures. Section~\ref{sec:prelim} states the conformal
framework and proves the two coverage results we rely on.
Section~\ref{sec:methods} describes the data and the experimental protocol.
Section~\ref{sec:results} presents the results, and Section~\ref{sec:discussion}
discusses their implications and limitations.

%=============================================================
\section{Related Work}\label{sec:related}

\subsection{Dropout prediction and educational data mining}

Institutional research on student attrition long predates its machine-learning
treatment. Tinto's integration model~\cite{tinto1975} framed departure as the
outcome of a student's academic and social integration, and supplied the
variables---prior attainment, financial circumstance, first-year
performance---that still dominate feature sets today. The move to predictive
modelling began with the educational data mining programme surveyed by Baker
and Yacef~\cite{baker2009} and updated by Romero and
Ventura~\cite{romero2020}; early applications to attrition, such as Kotsiantis
et al.~\cite{kotsiantis2003} on distance learning, established the template of
supervised classification from administrative records that the field has
followed since.

Work on the Portalegre cohort used here begins with Martins et
al.~\cite{martins2021}, who assembled the records and reported an initial
classification study, and with the data descriptor of Realinho et
al.~\cite{realinho2022} that placed the data in the public domain. Recent
studies on this cohort are methodologically careful. Gu and Zhang~\cite{gu2026}
compare eight classifiers---gradient boosting, LightGBM~\cite{ke2017},
XGBoost~\cite{chen2016}, CatBoost, random forests~\cite{breiman2001}, support
vector machines, AdaBoost and decision trees---and pair the best of them with
SHAP attributions~\cite{lundberg2017}, reporting a macro-AUC of 0.891 with a
confidence interval. Guam\'an-Hidalgo and Enciso~\cite{guaman2026} address the
three-class formulation directly, applying SMOTE~\cite{chawla2002} to the
well-known difficulty of learning from imbalanced data~\cite{he2009} and
clustering students into risk profiles.

What unites this literature is the level at which validity is assessed.
Whether the reported quantity is accuracy, macro-F1, AUC, or a confidence
interval on one of them, it summarises the model's behaviour over a cohort.
None of these quantities attaches a guarantee to the prediction made about a
particular student.

\subsection{Conformal prediction}

Conformal prediction~\cite{vovk2005,shafer2008} supplies exactly such a
guarantee. Given any fitted model and a held-out calibration sample, it returns
for each new instance a \emph{set} of labels containing the truth with a
user-specified probability, requiring only that calibration and test data be
exchangeable. The split, or inductive, variant introduced by Papadopoulos et
al.~\cite{papadopoulos2002} makes this practical by decoupling calibration from
model fitting, at the cost of one data split, and is widely used in applications; Lei et al.~\cite{lei2018} and the tutorial of
Angelopoulos and Bates~\cite{angelopoulos2023} give modern treatments.

The choice of nonconformity score determines the shape of the resulting sets.
The least-ambiguous-set score of Sadinle et al.~\cite{sadinle2019} minimises
average set size subject to the coverage constraint, while the adaptive
prediction sets of Romano et al.~\cite{romano2020} accumulate probability mass
in descending order and yield sets that adapt to instance difficulty. Both are
used here.

Applications now span domains in which a confident but incorrect prediction is
costly, and the framework has been extended to settings where exchangeability
fails only in structured ways, such as covariate shift~\cite{tibshirani2019}.
In education the framework is newer. Elmasry et al.~\cite{elmasry2026} apply
split conformal prediction to the present cohort within a broader trustworthy
modelling pipeline and report empirical coverage close to nominal. Their
analysis, like the applied conformal literature generally, reports that
coverage marginally.

\subsection{Conditional validity}

That distinction is the subject of established theory. In finite-label
classification, distribution-free object-conditional validity at non-atomic
covariates implies strong restrictions on label inclusion, often yielding
uninformative prediction sets in the relevant impossibility formulation \cite{vovk2012}. The infinite-length result concerns
regression intervals \cite{leiwasserman2014}; Foygel Barber et
al.~\cite{barber2021} analyse relaxed conditional coverage notions.
Mondrian conformal prediction calibrates within a finite taxonomy and gives
label-conditional validity when the taxonomy is the candidate outcome label
\cite{vovk2003mondrian}. This remedy to class-dependent miscoverage is not
novel: see the earlier class-conditional bias analysis of
L\"ofstr\"om et al.~\cite{lofstrom2015} and the class-imbalance discussion of
Ding et al.~\cite{ding2023}. A recent anomaly-detection audit by
Althani~\cite{althani2026} likewise studies class-specific coverage
under extreme imbalance; the present study differs in its three-category
educational outcome and measured feature-availability trade-offs.
Relative to the class-specific coverage audits discussed by Althani~\cite{althani2026},
we additionally examine five observed student attributes, their intersections
with realised outcomes, and the accompanying calibration-cell and set-size costs.
This is a dataset-specific empirical extension, not a new validity theorem.



\subsection{Selective prediction}

A prediction set with more than one label may be routed for additional
review, but set size alone does not guarantee the error rate among selected
singletons. Selective-classification risk control is a separate task:
Geifman and El-Yaniv~\cite{geifman2017,geifman2019} study risk-controlled
selection, and maximum model probability is a simple reference ranking.
Accordingly, we compare the \emph{delivered} Mondrian singleton label with
top-1 decisions on the same students and with a maximum-probability method
selecting an equal number of test cases. No guarantee is inferred from
$|C(x)|=1$.

%=============================================================
\section{Preliminaries}\label{sec:prelim}

We use standard split and Mondrian conformal prediction
\cite{vovk2005,papadopoulos2002,vovk2003mondrian}. The definitions and
finite-sample validity statements needed for the audit are given here;
proofs appear in Appendix~\ref{app:proofs}.

\subsection{Setting and notation}

Let $(X, Y)$ take values in $\mathcal{X} \times \Yset$, where $\mathcal{X}$ is
the space of student records and
\[
\Yset = \{\textsf{Dropout},\; \textsf{Enrolled},\; \textsf{Graduate}\}.
\]
Let $\pihat(y \mid x) \in [0,1]$ denote the score a fitted classifier assigns
to label $y$ for record $x$, with $\sum_{y \in \Yset} \pihat(y \mid x) = 1$. The
classifier is treated as fixed and arbitrary: no assumption is made about how it
was obtained, and it is not modified by any of what follows.

A \emph{set predictor} is a map $C : \mathcal{X} \to 2^{\Yset}$. We seek a
$C$ for which the true label lies in $C(X)$ with a prescribed probability.

\begin{definition}[Exchangeability; \cite{vovk2005}]
Random variables $Z_1,\dots,Z_m$ are \emph{exchangeable} if for every
permutation $\sigma$ of $\{1,\dots,m\}$ the joint law of
$(Z_{\sigma(1)},\dots,Z_{\sigma(m)})$ equals that of $(Z_1,\dots,Z_m)$.
\end{definition}

Exchangeability is strictly weaker than independence and is the only
distributional assumption used below. In particular, nothing is assumed about
the form of the law of $(X,Y)$, and no asymptotic regime is invoked.

\begin{definition}[Nonconformity score]
A \emph{nonconformity score} is a measurable map
$s : \mathcal{X} \times \Yset \to \mathbb{R}$ for which larger values indicate
that the pair $(x,y)$ conforms less well to the fitted model. Throughout we take
the \emph{least-ambiguous-set} (LAC) score of Sadinle et al.~\cite{sadinle2019},
\begin{equation}
s(x, y) \;=\; 1 - \pihat(y \mid x) ,
\label{eq:lac}
\end{equation}
and consider in addition the adaptive prediction sets (APS) score of Romano et
al.~\cite{romano2020} in Section~\ref{sec:robustness}.
\end{definition}

\subsection{The quantile lemma}

The entire construction rests on one elementary fact about the rank of an
exchangeable variable.

\begin{lemma}[Quantile lemma; \cite{vovk2005}, see also \cite{lei2018,angelopoulos2023}]\label{lem:quantile}
Let $Z_1, \dots, Z_n, Z_{n+1}$ be exchangeable real-valued random variables,
let $\alpha \in (0,1)$, and put $k = \lceil (n+1)(1-\alpha) \rceil$. If
$k \le n$, then
\begin{equation}
\Prob\big( Z_{n+1} \le Z_{(k)} \big) \;\ge\; 1-\alpha ,
\end{equation}
where $Z_{(1)} \le \dots \le Z_{(n)}$ denote the order statistics of
$Z_1,\dots,Z_n$. If in addition the $Z_i$ are almost surely distinct, then
$\Prob( Z_{n+1} \le Z_{(k)} ) \le 1-\alpha + \tfrac{1}{n+1}$.
\end{lemma}

\noindent The proof is provided in Appendix~\ref{app:proofs}.

\begin{remark}
The $(n+1)$ quantile correction preserves finite-sample validity:
for $n=885$ and $\alpha=0.10$, the correct order-statistic index is
798 rather than 797. The lemma's upper bound requires distinct scores.
\end{remark}

\subsection{Split conformal prediction and marginal validity}

Let $\mathcal{D}_{\mathrm{cal}} = \{(X_i, Y_i)\}_{i=1}^{n}$ be a calibration
sample disjoint from the data used to fit $\pihat$, and write
$s_i = s(X_i, Y_i)$. Define
\begin{equation}
\qhat \;=\; \begin{cases}s_{(k)},&k\le n,\\+\infty,&k>n,\end{cases}\qquad k=\big\lceil(n+1)(1-\alpha)\big\rceil,
\label{eq:qhat}
\end{equation}
and the split conformal set
\begin{equation}
C(x) \;=\; \big\{\, y \in \Yset \;:\; s(x,y) \le \qhat \,\big\}
\;=\; \big\{\, y \;:\; \pihat(y \mid x) \ge 1 - \qhat \,\big\} .
\label{eq:cset}
\end{equation}

\begin{theorem}[Marginal coverage; \cite{papadopoulos2002,vovk2005,lei2018}]\label{thm:marginal}
Suppose $(X_1,Y_1), \dots, (X_n, Y_n), (X_{n+1}, Y_{n+1})$ are exchangeable and
$\pihat$ is independent of them. Then, with $C$ as in \eqref{eq:cset},
\begin{equation}
\Prob\big( Y_{n+1} \in C(X_{n+1}) \big) \;\ge\; 1-\alpha .
\end{equation}
\end{theorem}

\noindent The proof is provided in Appendix~\ref{app:proofs}.

Theorem~\ref{thm:marginal} is distribution-free and finite-sample, and it is the
guarantee that split conformal prediction is normally reported as providing. It
is also, as stated, an average over the joint law of $(X,Y)$: it constrains no
subpopulation.

\subsection{Mondrian conformal prediction and class-conditional validity}

\begin{definition}[Taxonomy; \cite{vovk2003mondrian}]
A \emph{taxonomy} is a measurable map
$\kappa : \mathcal{X} \times \Yset \to \{1, \dots, K\}$ assigning each pair to
one of $K$ \emph{taxons}. We take the label taxonomy $\kappa(x,y) = y$, so that
the taxons are the outcome classes themselves.
\end{definition}

For each $c \in \Yset$ let $I_c = \{ i \le n : Y_i = c \}$, $n_c = |I_c|$, and
let $s^{(c)}_{(1)} \le \dots \le s^{(c)}_{(n_c)}$ be the ordered scores of the
calibration points in $I_c$. Put
\begin{equation}
\qhat_c \;=\; \begin{cases}s^{(c)}_{(k_c)},&k_c\le n_c,\\+\infty,&k_c>n_c,\end{cases}\qquad
k_c=\big\lceil(n_c+1)(1-\alpha)\big\rceil,
\label{eq:qhatc}
\end{equation}
and define the Mondrian conformal set
\begin{equation}
C^{\mathrm{M}}(x) \;=\; \big\{\, y \in \Yset \;:\; s(x,y) \le \qhat_y \,\big\} .
\label{eq:mset}
\end{equation}

\begin{theorem}[Class-conditional coverage; \cite{vovk2003mondrian,vovk2012}]\label{thm:conditional}
Suppose the fitted model is independent of calibration and test observations
and the observations are identically distributed, or are exchangeable
within each label stratum. For each $c\in\Yset$ with positive class
probability (including small or empty calibration strata under the
$+\infty$ convention),
\begin{equation}
\Prob\big( Y_{n+1} \in C^{\mathrm{M}}(X_{n+1}) \,\big|\, Y_{n+1} = c \big)
\;\ge\; 1-\alpha .
\end{equation}
\end{theorem}

\noindent The proof is provided in Appendix~\ref{app:proofs}.

\begin{remark}
Although the taxonomy is defined through the true label, $C^{\mathrm{M}}(x)$ is
computable without knowledge of $Y$: each candidate label $y$ is admitted or
rejected against its own threshold $\qhat_y$ in \eqref{eq:mset}. For a test
student whose realised outcome is $c$, admission is governed by $\qhat_c$, which
is what yields Theorem~\ref{thm:conditional}.
\end{remark}

\begin{remark}[Conditional coverage limits]\label{rem:impossible}
Coverage conditional on a prespecified finite label or group taxonomy
does not imply coverage conditional on every covariate value;
see~\cite{vovk2012,barber2021}. The infinite-length impossibility
result for regression intervals~\cite{leiwasserman2014} is not a
three-class set-size statement.
\end{remark}

%=============================================================
\section{Materials and Methods}\label{sec:methods}

\subsection{Data}

We use the publicly available \emph{Predict Students' Dropout and Academic
Success} dataset (UCI Machine Learning Repository, ID 697), collected at the
Polytechnic Institute of Portalegre and described by Realinho et
al.~\cite{realinho2022}. It records 4{,}424 students enrolled in 17
undergraduate programmes between the 2008/2009 and 2018/2019 academic years,
each described by 36 attributes covering demographic background, socio-economic
context, prior academic record, enrolment conditions, first- and
second-semester curricular performance, and three macroeconomic indicators.

The target is three-valued and recorded at the end of the normal duration of the
programme. Table~\ref{tab:data} summarises the dataset; it contains no missing
values and no duplicate records, as we verified prior to analysis.

\begin{table}[htbp]
\centering
\caption{Dataset description. UCI ID 697; no missing values, no duplicate rows.}
\label{tab:data}
\begin{tabular}{lrr}
\toprule
& \textbf{Count} & \textbf{Share} \\
\midrule
Records (students)      & 4{,}424 & \\
Attributes              & 36      & \\
Outcome classes         & 3       & \\
\midrule
\quad \textsf{Dropout}  & 1{,}421 & 32.12\% \\
\quad \textsf{Enrolled} &    794  & 17.95\% \\
\quad \textsf{Graduate} & 2{,}209 & 49.93\% \\
\bottomrule
\end{tabular}
\end{table}

Categorical attributes are supplied in the source data as integer codes and were
used in that form, without further encoding, for the tree-based models; inputs
to logistic regression were standardised. The primary baseline uses no hyperparameter tuning,
feature selection, or resampling; a separate training-only tuning sensitivity
analysis is reported below. This is deliberate: the object of study is the
behaviour of the conformal wrapper rather than the predictive ceiling of the
underlying classifier, and changes to preprocessing may alter class-specific score
distributions and conformal set sizes. We therefore interpret each
modelling comparison within its own stated experimental protocol.

\subsection{Data partitioning}

Split conformal prediction requires a calibration sample disjoint from the
training sample. We therefore partition the data three ways rather than two:
training (60\%, $n = 2{,}654$), calibration (20\%, $n = 885$) and test (20\%,
$n = 885$), stratified by outcome, so that the calibration and test parts each
contain 284 \textsf{Dropout}, 159 \textsf{Enrolled} and 442 \textsf{Graduate}
students.

The disjointness requirement is not a matter of good practice but a hypothesis
of Theorem~\ref{thm:marginal}: if calibration data influence the fitted model,
the calibration scores are no longer exchangeable with those of a future
student.

Earlier supplementary ten-partition sensitivity analyses used random seeds
$42$, $1$, $7$, $13$, $21$, $55$, $77$, $99$, $123$, and $2024$.
The primary coverage, robustness, feature, group, clustering, calibration-size, and selection audits
use 100 partitions (seeds 1000--1099). Across-partition standard deviations
and Monte Carlo standard errors describe split variability, not uncertainty
across independent student cohorts. Each partition uses two stratified
\texttt{train\_test\_split} calls with \texttt{random\_state=seed},
first allocating 60\% to training and then dividing the remaining 40\%
equally between calibration and test; the RF uses the same seed. Seed 42
is explicitly labelled wherever used.

\subsection{Base classifiers}

We evaluate three classifiers as the underlying point predictor: multinomial
logistic regression, a random forest~\cite{breiman2001} with 500 trees, and
histogram-based gradient boosting, all as implemented in
scikit-learn~\cite{pedregosa2011} and fitted on the training partition only. The
random forest serves as the base model throughout the conformal experiments; it
is a standard, reproducible reference model; its use does not imply
that it outperforms all alternative classifiers. Section 5.4 reports the alternative-model robustness experiments.
For the completed 100-partition robustness audit, logistic regression uses a training-only
\texttt{StandardScaler} pipeline with \texttt{LogisticRegression}
(\texttt{max\_iter=2000}). Histogram gradient boosting uses
\texttt{HistGradientBoostingClassifier} with the seed set and other
parameters at their defaults; the random forest uses 500 trees and the
partition seed. The seed-42 point-predictor comparison is descriptive and should not be
conflated with these completed 100-partition coverage experiments.

\subsection{Evaluation}

For a test sample of size $m$ we report empirical coverage, overall and
conditional on the realised class; the average prediction-set size
$m^{-1}\sum_i |C(X_i)|$ together with the full distribution of $|C(X_i)|$; and
\emph{singleton accuracy}, the proportion of correct labels among students whose
set contains exactly one.

Coverage and efficiency are always reported together. Coverage alone is not a
meaningful criterion, since the trivial predictor $C(x) = \Yset$ attains perfect
coverage at zero information; the question is what set size a given coverage
level costs.

Across repeated partitions we report the sample standard deviation of empirical
coverage and set size. Where the partition-average standard error is shown,
it is a Monte Carlo descriptive quantity (SD divided by the square root of
the number of partitions), not a sampling standard error for independent
student cohorts. The 100 partitions overlap substantially; no population-level
confidence statement is inferred from their apparent replication. Illustrative
seed-42 bootstrap and exact-binomial calculations, where retained, are
single-partition diagnostics rather than tests of conformal validity.

\subsection{Robustness protocol}\label{sec:robustness}

To assess model, score, calibration, and feature sensitivity, we evaluated
nine configurations over the same 100 stratified partitions (seeds 1000--1099).
The standard experiments use 60\% training, 20\% calibration, and 20\% test;
the doubled-calibration experiment instead uses 40\% training, 40\% calibration,
and 20\% test. Configurations include the baseline 500-tree random forest,
$\alpha=0.05$, randomized APS, histogram gradient boosting, logistic
regression, balanced class weights, and exclusion of six second-semester or
all twelve semester features (30 and 24 remaining predictors respectively).
For randomized APS, labels are ranked by decreasing fitted probability,
with the deterministic NumPy argsort tie convention. For each calibration
or test student, one $U\sim\mathrm{Uniform}(0,1)$ is shared across all
candidate labels, independently drawn in row order from
\texttt{numpy.random.default\_rng(seed)}. The candidate score is
$A(x,c)=\sum_{j:\,\mathrm{rank}(j)<\mathrm{rank}(c)}\hat\pi(j\mid x)
+\hat\pi(c\mid x)(1-U)$. Calibration uses the true-label score;
prediction includes each candidate whose score is at most the finite-sample
corrected threshold, including ties. These are label-wise randomized sets,
not necessarily threshold-crossing prefix sets. The original ten-seed
experiments are retained only where specifically identified as supplementary
or illustrative diagnostics.

\subsection{Reproducibility}

Analyses were conducted in local Jupyter using Python 3.13.9,
NumPy 2.3.5, pandas 2.3.3, and scikit-learn 1.7.2
\cite{pedregosa2011}. Primary analyses used seeds 1000--1099 and two
stratified splits per seed (60/20/20); separate supplementary
sensitivity analyses used their stated seeds. Source data are publicly
available from UCI 697. The analysis notebook, per-partition results,
and recorded split indices form the reproducibility materials. The
final repository and archival DOI will be inserted after public deposit.

%=============================================================
\section{Results}\label{sec:results}

\subsection{Primary base-model performance}

For the 100-partition reference random forest, mean test accuracy was
77.495\% and mean macro-F1 was 0.69267. All primary coverage comparisons
use seeds 1000--1099 and the specified 60/20/20 partition protocol.
Additional single-partition point-prediction results are retained in
Supplementary Section~S1 rather than treated as primary evidence.

\subsection{Marginal coverage is attained}

With $n=885$ calibration observations and $\alpha=0.10$, the finite-sample
quantile uses index $k=\lceil 886\times0.90\rceil=798$.
Across the 100 primary partitions, split CP attained mean marginal coverage
of 90.252\% (SD 1.503 percentage points) and mean prediction-set size
1.427 labels. These are repeated-partition averages, not a fixed-student
coverage guarantee.

\subsection{Class-conditional coverage is not}

Decomposing the same sets by realised outcome reveals substantial
class-specific differences (Table~\ref{tab:coverage} and Figure~\ref{fig:coverage}).

\begin{table}[htbp]
\centering\small
\caption{Primary class-wise coverage audit over 100 stratified partitions
(seeds 1000--1099; $\alpha=0.10$). Mean coverage and across-partition
standard deviation (SD) are percentages and percentage points, respectively;
SE is the descriptive Monte Carlo standard error of the partition mean.
``Below'' counts partitions with empirical coverage below the 90\% target,
not failures of the finite-sample conformal guarantee.}
\label{tab:coverage}
\begin{tabular}{llrrrr}
\toprule
Method & Class & Mean (\%) & SD (pp) & SE (pp) & Below/100 \\
\midrule
Split CP & Overall & 90.25 & 1.50 & 0.15 & 44 \\
 & \textsf{Dropout} & 88.04 & 2.09 & 0.21 & 82 \\
 & \textsf{Enrolled} & \textbf{71.73} & 5.61 & 0.56 & \textbf{100} \\
 & \textsf{Graduate} & 98.34 & 0.65 & 0.07 & 0 \\
\midrule
Mondrian CP & Overall & 90.40 & 1.34 & 0.13 & 37 \\
 & \textsf{Dropout} & 90.59 & 2.27 & 0.23 & 39 \\
 & \textsf{Enrolled} & \textbf{90.27} & 3.29 & 0.33 & 44 \\
 & \textsf{Graduate} & 90.33 & 2.00 & 0.20 & 40 \\
\bottomrule
\end{tabular}
\end{table}



\begin{figure}[htbp]
\centering
\begin{tikzpicture}
\begin{axis}[
  width=0.87\textwidth, height=6.0cm,
  ybar, bar width=11pt,
  ylabel={empirical coverage (\%)},
  xmin=-0.55, xmax=2.55,
  xtick={0,1,2}, xticklabels={Dropout,Enrolled,Graduate},
  ymin=50, ymax=106,
  ytick={50,60,70,80,90,100},
  grid=major, grid style={fdim!25},
  axis lines=left, tick align=outside,
  legend style={draw=none, at={(0.5,1.12)}, anchor=south, legend columns=2,
                font=\footnotesize, /tikz/every even column/.append style={column sep=10pt}},
  legend cell align=left,
  label style={font=\small}, tick label style={font=\footnotesize},
  error bars/y dir=both, error bars/y explicit,
  error bars/error bar style={black, thin},
  clip=false,
]
\addplot[fill=fgrey, draw=fgrey, bar shift=-7pt] coordinates {
  (0,88.04) +- (0,2.09)
  (1,71.73) +- (0,5.61)
  (2,98.34) +- (0,0.65)};
\addlegendentry{split CP}
\addplot[fill=fblue, draw=fblue, bar shift=7pt] coordinates {
  (0,90.59) +- (0,2.27)
  (1,90.27) +- (0,3.29)
  (2,90.33) +- (0,2.00)};
\addlegendentry{Mondrian CP}

% nominal level -- drawn, not plotted, so it is not treated as a third bar series
\draw[black, dashed, thin]
  (axis cs:-0.55,90) -- (axis cs:2.55,90);
% The dashed line is identified in the caption, avoiding overlap with bars.
\end{axis}
\end{tikzpicture}
\caption{Class-conditional coverage under split and Mondrian conformal
prediction. Bars are means over 100 random partitions (seeds 1000--1099); error bars are one
standard deviation; the dashed line is the nominal level. Random forest
(500 trees), stratified 60/20/20 split, $\alpha = 0.10$, LAC score.}
\label{fig:coverage}
\end{figure}

Over 100 partitions, standard split CP covers only 71.73\% of \textsf{Enrolled}
students on average, a shortfall of 18.27 percentage points, and every
partition falls below 90\%. The marginal guarantee is satisfied not because it holds across the
population but because the largest and easiest class is over-covered by enough
to compensate for the two harder ones. Split conformal prediction meets its
promise on average by over-serving \textsf{Graduate} and under-serving both
\textsf{Dropout} and \textsf{Enrolled}.

A supplementary seed-42 score-distribution illustration provides a
single-partition diagnostic; the quantitative conclusions here use the
100-partition empirical coverage estimates above.

The test partition contains approximately 159 \textsf{Enrolled} students, so
individual-split coverage estimates are noisy. The repeated-split mean and SD
above characterize the empirical sensitivity to repartitioning. They do not
establish a conditional guarantee for a fixed student or justify treating the
100 overlapping splits as independent populations.

\subsection{The failure is not an artefact of design choices}

Table~\ref{tab:robust} presents the 100-partition robustness comparison.
Despite approximately nominal marginal coverage, each reported split-CP
configuration has an average \textsf{Enrolled} coverage below its nominal target;
all 100 empirical partitions fall below the applicable target in each
configuration. Randomized APS increases mean \textsf{Enrolled} coverage
relative to LAC but does not eliminate the shortfall.

\begin{table}[htbp]
\centering\small
\caption{Robustness of standard split CP across 100 partitions (seeds
1000--1099). Coverage is a percentage; SD is the across-partition standard
deviation in percentage points. Failures count test partitions with empirical
\textsf{Enrolled} coverage below the applicable 90\% or 95\% target.}
\label{tab:robust}
\begin{tabular}{lrrrr}
\toprule
\textbf{Variation} & \textbf{Enrolled} & \textbf{SD} & \textbf{Overall} & \textbf{Fails}\\
\midrule
RF LAC, $\alpha=0.10$ & 71.730 & 5.606 & 90.252 & 100/100\\
RF LAC, $\alpha=0.05$ (95\% target) & 86.679 & 3.614 & 95.241 & 100/100\\
RF randomized APS & 78.836 & 4.310 & 90.173 & 100/100\\
RF, 40\% calibration & 71.126 & 4.358 & 90.060 & 100/100\\
Histogram gradient boosting & 72.723 & 4.871 & 90.160 & 100/100\\
Logistic regression & 70.113 & 5.184 & 90.120 & 100/100\\
\bottomrule
\end{tabular}
\end{table}

At $\alpha=0.05$, the 86.679\% Enrolled mean remains below the 95\%
target. The randomized APS score, defined in Section~\ref{sec:robustness},
gives 78.836\% Enrolled coverage compared with 71.730\% under LAC;
the altered score changes the coverage pattern without implying classwise
validity for global calibration.

A deterministic cumulative-score sensitivity analysis using ten partitions
is reported separately in Supplementary Section~S4.

Increasing the calibration share cannot be interpreted as an isolated
calibration-size intervention. Raising calibration from 20\% to 40\% increases its count
from 885 to 1,770 and concurrently reduces training from 2,654 to 1,769
observations, with test size held at 885. The resulting 71.126\% mean Enrolled
coverage cannot identify the separate effect of calibration size. A single marginal quantile cannot simultaneously
serve three score distributions at different locations, however precisely it is
estimated.

\subsection{Mondrian conformal prediction restores class-conditional coverage}

Label-Mondrian CP uses a separate calibration threshold for each candidate
class. The 100-partition comparisons below quantify its coverage and
set-size effects; illustrative seed-42 thresholds are provided in the
supplement.



Table~\ref{tab:coverage} presents the matched 100-partition
comparison. The increase for \textsf{Enrolled} is 18.540 percentage points,
while \textsf{Graduate} coverage decreases by 8.011 points from its large
split-CP overcoverage. The between-class coverage spread declines from
26.607 to approximately 0.318 percentage points. Empirical test coverage
nevertheless falls below 90\% in some individual partitions: under Mondrian,
39 Dropout, 44 Enrolled, and 40 Graduate partitions, and 37 overall, fall
below the target. This does not contradict the repeated-sampling guarantee.

The resulting set-size trade-off is summarised below; a seed-42
set-size histogram is provided in Supplementary Section~S3.



\paragraph{Across-partition variability and prediction-set efficiency.}
The consolidated results in Table~\ref{tab:coverage} provide means, SD,
SE, and below-target counts for both methods on the same 100 partitions.
Average set size increases from 1.427 (split CP) to 1.606
(label-Mondrian CP), a difference of 0.179 labels per student.

\subsection{Class weighting and earlier-feature experiments}\label{sec:featurechecks}

The additional 100-partition results are in Table~\ref{tab:features}.
The weighted RF uses \texttt{class\_weight=balanced}, not a distinct balanced
random-forest resampling algorithm. The earlier-feature models retain 30
non-second-semester predictors or 24 predictors excluding both semesters.
The latter is a feature-availability proxy, not longitudinal validation.

\begin{table}[htbp]
\centering\small
\caption{Additional model/feature robustness checks across 100 partitions
(seeds 1000--1099). SD and below-nominal counts refer to Enrolled coverage;
Mondrian set size is averaged over all test students.}
\label{tab:features}
\begin{tabular}{lrrrrr}
\toprule
Variation & Split Enrolled & SD & Below 90\% & Mondrian Enrolled & Mondrian size\\
\midrule
RF balanced & 70.610\% & 5.012 & 100/100 & 90.346\% & 1.600\\
No second semester (30) & 73.409\% & 4.611 & 100/100 & 90.308\% & 1.724\\
No semester features (24) & 69.503\% & 5.240 & 100/100 & 90.440\% & 2.126\\
\bottomrule
\end{tabular}
\end{table}

The 24-predictor configuration has 69.503\% mean Enrolled split-CP coverage
but 90.440\% under label-Mondrian calibration, with larger mean sets
(2.126 labels). This illustrates the validity--efficiency trade-off when
post-enrolment academic features are unavailable.

\subsection{Demographic and group-by-label coverage audit}\label{sec:groups}

We examined five observed binary attributes (Gender, Scholarship holder,
Debtor, Displaced, International). Within each of 100 stratified partitions
(seeds 1000--1099), three methods used the same fitted 500-tree random forest:
standard split CP with a single threshold, group-Mondrian CP with a threshold
for each observed attribute value, and group-by-label Mondrian CP with a
threshold for each attribute--candidate-label pair. All thresholds used the
finite-sample index $\lceil(n+1)(1-\alpha)\rceil$ within their corresponding
calibration taxon, assigning $+\infty$ when that index exceeded the available
calibration count. The latter choice is conservative but may result in very
large prediction sets. Attributes were examined separately, not as a joint
five-variable intersection. The fitted predictor was identical across methods.

\begin{table}[htbp]
\centering\small
\caption{Coverage (\%) and mean calibration-cell counts over 100 partitions
for realised \textsf{Enrolled} students. Counts are mean labelled
calibration observations in the displayed attribute--outcome cell;
values are descriptive and no multiple-testing inference is made. For International = 1 among Enrolled test students, coverage uses 99 evaluable partitions (seed 1037 has no such test case); calibration counts average all 100 partitions.}
\label{tab:groupaudit}
\begin{tabular}{llrrr}
\toprule Attribute & Value & Mean cal. $n$ & Split CP & Group-by-label CP\\
\midrule
Gender & 0 & 96.79 & 71.47 & 90.77\\
Gender & 1 & 62.21 & 72.15 & 91.40\\
Scholarship holder & 0 & 132.34 & 73.92 & 90.67\\
Scholarship holder & 1 & 26.66 & 60.79 & 91.46\\
Debtor & 0 & 140.81 & 71.94 & 90.56\\
Debtor & 1 & 18.19 & 70.00 & 92.83\\
Displaced & 0 & 72.13 & 73.16 & 91.01\\
Displaced & 1 & 86.87 & 70.59 & 90.66\\
International & 0 & 154.38 & 71.39 & 90.78\\
International & 1 & 4.62 & 82.96 & 99.80\\
\bottomrule
\end{tabular}
\end{table}

\begin{table}[htbp]
\centering\small
\caption{Mean empirical coverage (\%) over 100 partitions within each observed attribute group, pooling realised labels. Group-Mondrian calibration is compared with both global split CP and group-by-label calibration; these group-level averages do not establish validity for every group--label cell.}
\label{tab:groupoverall}
\begin{tabular}{llrrr}
\toprule Attribute & Value & Split CP & Group-Mondrian CP & Group-by-label CP\\
\midrule
Gender & 0 & 90.52 & 90.24 & 90.61\\
Gender & 1 & 89.75 & 90.49 & 90.84\\
Scholarship holder & 0 & 89.99 & 90.25 & 90.46\\
Scholarship holder & 1 & 91.04 & 90.24 & 90.80\\
Debtor & 0 & 90.17 & 90.27 & 90.45\\
Debtor & 1 & 90.91 & 91.24 & 92.03\\
Displaced & 0 & 91.56 & 90.34 & 90.85\\
Displaced & 1 & 89.17 & 90.22 & 90.30\\
International & 0 & 90.22 & 90.25 & 90.51\\
International & 1 & 91.48 & 91.90 & 96.60\\
\bottomrule
\end{tabular}
\end{table}

The 30.67-percentage-point increase for \textsf{Enrolled} scholarship holders
(60.79\% to 91.46\%) is an empirical contrast, not evidence of improved
point prediction. Rare cells may yield nearly vacuous sets: among international
students, group-by-label Mondrian CP attains 96.60\% mean overall coverage
with mean set size 2.60, versus 91.48\% and size 1.56 under split CP.
The international-group calibration count averages only 21.95 observations
per partition across all outcomes. The per-partition calibration counts are retained in the computational
results file, with mean \textsf{Enrolled} cell counts reported in Table~\ref{tab:groupaudit}.
These subgroup contrasts are descriptive; no multiple-testing inference is made.
Table~\ref{tab:groupoverall} additionally documents the group-Mondrian
baseline, which does not by itself target label-conditional coverage within groups.

\subsection{Two-cluster score-distribution sensitivity analysis}\label{sec:cluster}

As a small-class exploratory comparison motivated by
Ding et al.~\cite{ding2023}, we used 25\% of each held-out calibration
partition for estimating five quantiles (.10, .25, .50, .75, .90) of true-label
LAC scores within each outcome and applied two-cluster $k$-means to the three
class-level quantile vectors. The remaining 75\% of calibration records formed
an independent quantile-calibration subset for the clustered, global, and
classwise procedures. The clustering method and the number of clusters are
an adaptation for three outcomes, \emph{not} an exact implementation of
Ding et al.'s complete many-class procedure and cannot support a claim of exact replication. The clustering is fixed before final
quantile calibration to prevent reuse of those scores for constructing strata.

\begin{table}[htbp]
\centering\small
\caption{100-partition comparison using the common 75\% calibration subset;
coverage in \%, with overall mean prediction-set size.}
\label{tab:cluster}
\begin{tabular}{lrrrrr}
\toprule Method & Overall & Dropout & Enrolled & Graduate & Size\\
\midrule
Split CP & 90.16 & 87.94 & 71.43 & 98.33 & 1.423\\
Two-cluster CP & 90.27 & 82.00 & 90.45 & 95.52 & 1.543\\
Classwise CP & 90.36 & 90.34 & 90.45 & 90.34 & 1.604\\
\bottomrule
\end{tabular}
\end{table}

The clustered variant maintains near-nominal marginal coverage but
under-covers \textsf{Dropout} on average (82.00\%, below 90\% on 100/100
partitions). Classwise calibration has higher mean set size but balances
coverage across the three outcomes. Cluster-conditional validity does not in
general imply label-conditional validity within a multi-label cluster.

\subsection{Fixed-training calibration-size ablation}\label{sec:fixedcal}

To isolate calibration sample size from model-fitting sample size, each of
100 partitions retained the same 60\% training subset (2,654 records)
and 20\% test subset (885 records) while using a stratified 25\%, 50\%,
or 100\% subsample of the original 885-record calibration pool. Each RF
was fitted only once per seed. The two smaller calibration subsets were
sampled separately and are not necessarily nested.

\begin{table}[htbp]
\centering\small
\caption{Fixed-training \textsf{Enrolled} coverage (\%) by calibration sample
size; 100-partition mean and standard deviation (SD), in percentage points.}
\label{tab:fixedcal}
\begin{tabular}{rrrrr}
\toprule $n_{\mathrm{cal}}$ & Split mean & Split SD & Mondrian mean & Mondrian SD\\
\midrule
221 & 73.09 & 7.17 & 90.20 & 5.17\\
442 & 71.92 & 6.15 & 90.52 & 4.07\\
885 & 71.73 & 5.61 & 90.27 & 3.29\\
\bottomrule
\end{tabular}
\end{table}

Increasing calibration count reduces between-partition variability but
fails to remove the class-specific gap under a global threshold. These
results do not imply a monotonic change in mean coverage as calibration
count increases.

\subsection{Training-only tuning: supplementary sensitivity}

Separately, per-partition three-fold training-only hyperparameter tuning
was evaluated on ten predefined seeds. It did not improve split-CP
\textsf{Enrolled} coverage: the default and tuned mean coverages were
73.145\% and 72.830\%, respectively. The tuning experiment is not a
100-partition analysis; its full protocol, table, and across-seed
variability are reported in Supplementary Section~S5.

\subsection{Selective-decision comparison}\label{sec:triage}

The primary selection comparison uses 100 partitions and exactly matched
automatic-decision counts per partition. A separate seed-42 comparison is
provided in Supplementary Section~S6.

\begin{table}[htbp]
\centering
\caption{Selective prediction on 100 partitions (60/20/20; $\alpha=0.10$).
Methods select an exactly matched number of test students per partition.
SD is across partitions; differences in accuracy are percentage points.
Neither selected group has a conformal conditional error-rate guarantee.
The identical three-decimal means of the two Mondrian accuracy measures
are observed rounded sample averages; the measures differ within individual
partitions (for example, seed 42) and have different SDs. The paired difference has mean 0.0003
percentage points and SD 0.6095 percentage points, and is nonzero in 71 of
100 partitions.}
\label{tab:triage}
\begin{tabular}{lrr}
\toprule
Metric & Mean & SD\\
\midrule
Automatically decided share (\%) & 44.168 & 4.461\\
Mondrian delivered singleton accuracy (\%) & 90.718 & 2.108\\
Mondrian top-1 accuracy on selected (\%) & 90.718 & 1.988\\
Max-probability matched-selection accuracy (\%) & 93.322 & 1.433\\
Mondrian delivered minus same-student top-1 (pp) & 0.0003 & 0.6095\\
Max-probability minus Mondrian delivered (pp) & 2.6033 & 1.6180\\
Partitions with nonzero delivered--top-1 difference & 71/100 & ---\\
Partitions with nonzero max-probability difference & 98/100 & ---\\
Empty sets per test partition & 0.000 & 0.000\\
\bottomrule
\end{tabular}
\end{table}

At the standard $\alpha=0.10$, the non-selected students are the referred
group. The pooled referred-group composition over 100 test partitions is
32.12\% Dropout, 26.59\% Enrolled and 41.30\% Graduate for the
singleton rule, compared with 29.40\%, 28.76\% and 41.84\% under matched
maximum-probability selection. The denominator is the number \emph{referred}
under each rule, rather than the number belonging to each class. Both methods
yield 49,411 pooled referrals across 100 overlapping test partitions; these
are not unique individuals. Exact pooled referral counts are 15,870, 13,136,
and 20,405 (Mondrian) versus 14,529, 14,210, and 20,672 (matched
maximum-probability selection) for Dropout, Enrolled, and Graduate,
respectively. Per-partition checks verified that the three referred-class
counts sum to test size minus the number of automatically decided students
for both selection rules.

A single-partition sweep over miscoverage levels is reported separately
in Supplementary Section~S6; the primary selection findings follow from
the 100-partition paired comparisons in Table~\ref{tab:triage}.

%=============================================================
\section{Discussion}\label{sec:discussion}

\subsection{Summary}

In the 100-partition primary analysis, standard split CP has mean marginal
coverage 90.25\% but only 71.73\% for Enrolled, below nominal in every
partition. The robustness and feature-availability audits use the same 100 partitions.
The discrepancy reflects class-specific score distributions.
The 100-partition Mondrian class means are close to the nominal target,
with larger sets. The matched confidence baseline outperforms singleton
selection in accuracy, so the singleton rule should not be presented as a
superior risk-controlled triage method.

\subsection{Relation to prior work on this cohort}

Gu and Zhang~\cite{gu2026} report a confidence interval on an AUC. That
interval quantifies how precisely average discriminative ability has been
estimated; it is not, and is not intended to be, a statement about an individual
student. A confidence interval on a population parameter and a coverage
guarantee for a prediction set answer different questions, and the present work
supplies the second.

Guamán-Hidalgo and Enciso~\cite{guaman2026} consider three-class
outcome modelling; our contribution concerns the conditional interpretation of
prediction-set coverage rather than a head-to-head comparison of their
predictive performance.

Elmasry et al.~\cite{elmasry2026} report near-nominal marginal conformal
coverage on this benchmark. The present experiments show that such a marginal
summary can coexist with considerably lower coverage for a particular outcome
or demographic--outcome stratum. 

\subsection{Why marginal coverage is the wrong target here}

Marginal and label-conditional coverage are distinct requirements. A shared
threshold meets both only if each class-specific score CDF at that threshold
reaches the nominal coverage level. On these data the class-specific empirical coverages differ materially
at the global threshold. A single quantile applied to three distributions at different locations must
under-serve those on the right and over-serve those on the left, and increasing
the calibration sample refines the estimate of that quantile without altering
this fact---as the fixed-training calibration-subset experiment further illustrates.

Finite-taxonomy conditional validity rests on exchangeability within
the specified calibration groups; it differs from instance-conditional validity.

\subsection{Practical implications}

For this benchmark, practitioners should report both marginal and
class-conditional empirical coverage, using independent calibration data.
Mondrian's thresholds use \emph{labelled calibration data} separately for each
candidate class; they are not inferred autonomously from unlabelled students.
It can enlarge sets for harder groups, but set size does not itself guarantee
selective risk. Any operational referral workflow requires evaluation against
confidence-based ranking and other risk-control methods; here the simple
matched maximum-probability comparison was more accurate.

\noindent\fbox{\begin{minipage}{0.94\linewidth}
\textbf{Recommended reporting checklist for educational conformal prediction.}
Report (i) marginal coverage; (ii) coverage separately by outcome label;
(iii) coverage for prespecified student groups and group--outcome cells,
with calibration-cell counts; (iv) mean prediction-set size and its
distribution; and (v) automatic-decision performance against a matched,
equal-capacity confidence-based selection baseline.
\end{minipage}}

\subsection{Limitations}

\textbf{Exchangeability.} The random splits are disjoint within each partition, but this alone does
not establish exchangeability of observations in the underlying population.
The finite-sample guarantees require exchangeability of calibration and
test observations under the sampling model; random partitioning does not
validate that assumption for temporal deployment. It does not follow that a model calibrated on one cohort retains
its guarantee on a later one: the data span 2008/2009 to 2018/2019, over which
curricular, economic and admissions conditions change. We have not tested
temporal transfer, and a deployed system would require periodic recalibration.

\textbf{Single institution.} All students attend one Portuguese polytechnic.
Whether the same class-conditional pattern arises elsewhere is an empirical
question these data cannot answer, though the mechanism we identify---a score
distribution differing by class---is not specific to this institution.

\textbf{Calibration sample for the minority class.} The \textsf{Enrolled}
threshold is estimated from 159 calibration students. The classwise quantile is therefore estimated more coarsely than
those for the two larger classes. We separately report a fixed-training
calibration-size study to avoid interpreting the earlier confounded 40\%
calibration comparison as an isolated calibration intervention.

\textbf{Ambiguity of the target class.} \textsf{Enrolled} denotes students still
registered at the census point rather than a settled outcome, and is therefore
censored in a way the other two classes are not. Part of its difficulty may be
intrinsic to that definition. This does not affect the coverage results, which
concern whatever label the data records, but it bears on their interpretation.

\textbf{Modelling choices.} Categorical attributes are used in their supplied
integer coding. The primary baseline is untuned, but a separate training-only
tuning experiment is included. These choices
were deliberate and the phenomenon is robust across three model families, but a
more carefully engineered classifier would shift the absolute coverage figures,
if not the pattern.

\subsection{Future work}

Three extensions follow naturally. The group and group-by-label audits
performed here could be extended to jointly intersecting attributes and
external institutions, with careful attention to sparse strata. Temporal validation would require access to reliable cohort or entry-year
metadata, which are not included as an explicit variable in the supplied
36-feature model matrix. Such a study cannot be conducted from these
features alone. Finally, score
functions designed for class-conditional validity may recover the efficiency
lost here.

%=============================================================
\section{Conclusions}

This cohort-specific audit documents the known distinction between marginal
and label-conditional conformal coverage. Across 100 random partitions, a
random-forest split-CP model achieved 90.25\% mean marginal coverage but only
71.73\% Enrolled coverage (below 90\% on all 100 partitions). Mondrian CP
achieved 90.59\%, 90.27\% and 90.33\% mean coverage for the three outcome
classes, consistent with the class-conditional validity theorem, while
increasing mean prediction-set size from 1.427 to 1.606.

The robustness and earlier-feature studies show that undercoverage can persist
under different predictive settings. A fixed-training calibration-size ablation indicates that increasing
calibration count alone does not remove the global split-CP Enrolled gap.
Group-by-label conditioning can address additional intersectional gaps
at the cost of larger sets, while a two-cluster adaptation does not ensure
per-class coverage. A supplementary ten-seed training-only RF tuning analysis showed no
improvement in Enrolled split-CP coverage. Finally, automatically delivering Mondrian singletons was less accurate than
selecting an equal number of high-confidence predictions (90.72\% versus
93.32\% over 100 partitions). No singleton conditional-risk guarantee follows
from conformal coverage. For educational applications, class-wise auditing and
explicit coverage--efficiency--selection comparisons are more informative
than a single marginal coverage estimate.

%=============================================================
\section*{Data Availability Statement}

The dataset analysed in this study is publicly available from the UCI Machine
Learning Repository as ``Predict Students' Dropout and Academic Success''
(ID 697), at
\url{https://archive.ics.uci.edu/dataset/697/predict+students+dropout+and+academic+success},
under a Creative Commons Attribution 4.0 International (CC BY 4.0) licence. The
data were originally collected at the Polytechnic Institute of Portalegre,
Portugal, and are described by Realinho et al.~\cite{realinho2022}. The version
used here comprises 4{,}424 records with 36 attributes and a three-class target,
with no missing values. Accessed 8 September 2026.

\section*{Author Contributions}
Rimsha Bibi: conceptualization, methodology, software, formal analysis,
investigation, visualization, and writing---original draft.
Dr. M. Umair Shahzad: supervision, review and editing.
% Pre-submission: both authors must approve their exact CRediT roles.

\section*{Ethics Statement}
The study is a secondary analysis of publicly available,
de-identified student records; no students were recruited or contacted.
% Pre-submission: institutional policy must be checked before
% adding a claim that ethical approval was unnecessary or waived.

\section*{Funding}
No external grant funding was received for this research.
% Pre-submission: both authors must confirm funding declaration.

\section*{Code Availability}
The analysis code and reproducibility outputs are being prepared for
public archiving. The public repository URL and persistent Zenodo DOI
will be provided once the deposit is complete.
% REQUIRED BEFORE SUBMISSION: replace this text with actual verified links.

\section*{AI-use Declaration}
AI-assisted editorial and code-drafting tools were used during manuscript
and analysis preparation. The authors retain responsibility for the
accuracy, reproducibility and interpretation of the submitted work.
% Pre-submission: verify AI-use wording meets journal policy and both authors approve.

%=============================================================
\appendix
\section{Proofs of the conformal validity statements}\label{app:proofs}

\subsection{Quantile lemma (Lemma~\ref{lem:quantile})}
\begin{proof}
Introduce independent $U_1,\ldots,U_{n+1}\sim\mathrm{Uniform}(0,1)$,
independent of the exchangeable $Z_i$, and order pairs $(Z_i,U_i)$
lexicographically. Almost surely these pairs have distinct ranks and are
exchangeable, so the rank $R'$ of $(Z_{n+1},U_{n+1})$ among all $n+1$ pairs
is uniform on $\{1,\ldots,n+1\}$. If $R'\le k$, then
$Z_{n+1}\le Z_{(k)}$, including when original scores are tied. Hence
\[
\Prob\{Z_{n+1}\le Z_{(k)}\}\ge\Prob\{R'\le k\}
=\frac{k}{n+1}\ge1-\alpha.
\]
If the original $Z_i$ are almost surely distinct, the relevant ranks
coincide, giving equality $k/(n+1)$ and the claimed upper bound.


\end{proof}

\subsection{Marginal coverage (Theorem~\ref{thm:marginal})}
\begin{proof}
Since $s$ is a fixed measurable function of a pair and $\pihat$ is independent of
the sample, the scores $s_1, \dots, s_n, s_{n+1}$ inherit exchangeability from
the pairs. By construction $Y_{n+1} \in C(X_{n+1})$ if and only if
$s_{n+1} \le \qhat$. If $k\le n$, Lemma~\ref{lem:quantile} applies;
if $k>n$, the threshold is $+\infty$ and the candidate label is always
included. Both cases establish the claim.
\end{proof}

\subsection{Class-conditional coverage (Theorem~\ref{thm:conditional})}
\begin{proof}
Condition on the class-label sequence and the event $Y_{n+1}=c$.
Under the usual independent, identically distributed sampling of students
(and, more generally, appropriate within-taxonomy exchangeability), the
calibration observations with label $c$ and the test observation with
label $c$ are exchangeable, and therefore their scores
$\{s_i : i \in I_c\} \cup \{s_{n+1}\}$ are exchangeable. On this event,
$Y_{n+1} \in C^{\mathrm{M}}(X_{n+1})$ holds if and only if
$s_{n+1} \le \qhat_c$, since admission of the label $c$ is governed by
$\qhat_c$ alone. Lemma~\ref{lem:quantile} applied with $n_c$ in place of $n$
gives coverage at least $1-\alpha$ when $k_c\le n_c$. For $k_c>n_c$,
including an empty calibration stratum, the $+\infty$ threshold ensures
inclusion. Averaging over the conditioning variables preserves the bound.
\end{proof}

%=============================================================
\begin{thebibliography}{99}
\setlength{\itemsep}{2pt}

\bibitem{realinho2022}
Realinho, V.; Machado, J.; Baptista, L.; Martins, M.V.
Predicting Student Dropout and Academic Success.
\emph{Data} \textbf{2022}, \emph{7}, 146.
\url{https://doi.org/10.3390/data7110146}
\bibitem{gu2026}
Gu, H.; Zhang, Y.
Efficient and Interpretable Machine Learning for Student Academic Outcome
Prediction.
\emph{Mathematics} \textbf{2026}, \emph{14}, 626.
\url{https://doi.org/10.3390/math14040626}
\bibitem{guaman2026}
Guam\'an-Hidalgo, E.; Enciso, L.
Early Prediction of Student Dropout Through Machine Learning in Educational
Data Mining: A Comparative Analysis of Algorithms and Data Balancing
Techniques.
\emph{Applied Sciences} \textbf{2026}, \emph{16}, 8915.
\url{https://doi.org/10.3390/app16188915}
\bibitem{elmasry2026}
Elmasry, M.M.S.; Gafar, M.G.; Elsabagh, M.A.
Trustworthy Educational Risk Modeling with Calibrated Probabilities, Conformal
Uncertainty, Explainable AI, and Graph-Based Refinement.
\emph{Inventions} \textbf{2026}, \emph{11}, 65.
\url{https://doi.org/10.3390/inventions11030065}

% --- education and educational data mining ---
\bibitem{ding2023}
Ding, T.; Angelopoulos, A.N.; Bates, S.; Jordan, M.I.; Tibshirani, R.J.
Class-Conditional Conformal Prediction with Many Classes.
In \emph{Advances in Neural Information Processing Systems 36}, 2023.
\url{https://proceedings.neurips.cc/paper_files/paper/2023/hash/cb931eddd563f8d473c355518ce8601c-Abstract-Conference.html}
\bibitem{lofstrom2015}
L\"ofstr\"om, T.; Bostr\"om, H.; Linusson, H.; Johansson, U.
Bias reduction through conditional conformal prediction.
\emph{Intelligent Data Analysis} \textbf{2015}, \emph{19}(6), 1355--1375.
\url{https://doi.org/10.3233/IDA-150786}
\bibitem{tinto1975}
Tinto, V.
Dropout from Higher Education: A Theoretical Synthesis of Recent Research.
\emph{Review of Educational Research} \textbf{1975}, \emph{45}, 89--125.
\url{https://doi.org/10.3102/00346543045001089}
\bibitem{baker2009}
Baker, R.S.J.d.; Yacef, K.
The State of Educational Data Mining in 2009: A Review and Future Visions.
\emph{Journal of Educational Data Mining} \textbf{2009}, \emph{1}, 3--17.
\url{https://jedm.educationaldatamining.org/index.php/JEDM/article/view/8}
\bibitem{romero2020}
Romero, C.; Ventura, S.
Educational Data Mining and Learning Analytics: An Updated Survey.
\emph{WIREs Data Mining and Knowledge Discovery} \textbf{2020}, \emph{10},
e1355.
\url{https://doi.org/10.1002/widm.1355}
\bibitem{kotsiantis2003}
Kotsiantis, S.B.; Pierrakeas, C.J.; Pintelas, P.E.
Preventing Student Dropout in Distance Learning Using Machine Learning
Techniques.
In \emph{Knowledge-Based Intelligent Information and Engineering Systems};
Lecture Notes in Computer Science 2774; Springer: Berlin, 2003; pp.~267--274.
\url{https://doi.org/10.1007/978-3-540-45226-3_37}
\bibitem{martins2021}
Martins, M.V.; Tolledo, D.; Machado, J.; Baptista, L.M.T.; Realinho, V.
Early Prediction of Student's Performance in Higher Education: A Case Study.
In \emph{Trends and Applications in Information Systems and Technologies};
Advances in Intelligent Systems and Computing 1365; Springer: Cham, 2021;
pp.~166--175.
\url{https://doi.org/10.1007/978-3-030-72657-7_16}

% --- learning algorithms and interpretability ---
\bibitem{ke2017}
Ke, G.; Meng, Q.; Finley, T.; Wang, T.; Chen, W.; Ma, W.; Ye, Q.; Liu, T.-Y.
LightGBM: A Highly Efficient Gradient Boosting Decision Tree.
In \emph{Advances in Neural Information Processing Systems 30}, 2017;
pp.~3146--3154.
\url{https://papers.nips.cc/paper_files/paper/2017/hash/6449f44a102fde848669bdd9eb6b76fa-Abstract.html}
\bibitem{chen2016}
Chen, T.; Guestrin, C.
XGBoost: A Scalable Tree Boosting System.
In \emph{Proceedings of the 22nd ACM SIGKDD International Conference on
Knowledge Discovery and Data Mining}, 2016; pp.~785--794.
\url{https://doi.org/10.1145/2939672.2939785}
\bibitem{breiman2001}
Breiman, L.
Random Forests.
\emph{Machine Learning} \textbf{2001}, \emph{45}, 5--32.
\url{https://doi.org/10.1023/A:1010933404324}
\bibitem{lundberg2017}
Lundberg, S.M.; Lee, S.-I.
A Unified Approach to Interpreting Model Predictions.
In \emph{Advances in Neural Information Processing Systems 30}, 2017;
pp.~4765--4774.
\url{https://papers.nips.cc/paper_files/paper/2017/hash/8a20a8621978632d76c43dfd28b67767-Abstract.html}
\bibitem{chawla2002}
Chawla, N.V.; Bowyer, K.W.; Hall, L.O.; Kegelmeyer, W.P.
SMOTE: Synthetic Minority Over-sampling Technique.
\emph{Journal of Artificial Intelligence Research} \textbf{2002}, \emph{16},
321--357.
\url{https://doi.org/10.1613/jair.953}
\bibitem{he2009}
He, H.; Garcia, E.A.
Learning from Imbalanced Data.
\emph{IEEE Transactions on Knowledge and Data Engineering} \textbf{2009},
\emph{21}, 1263--1284.
\url{https://doi.org/10.1109/TKDE.2008.239}

% --- conformal prediction ---
\bibitem{vovk2005}
Vovk, V.; Gammerman, A.; Shafer, G.
\emph{Algorithmic Learning in a Random World}; Springer: New York, 2005.
\url{https://doi.org/10.1007/b106715}
\bibitem{shafer2008}
Shafer, G.; Vovk, V.
A Tutorial on Conformal Prediction.
\emph{Journal of Machine Learning Research} \textbf{2008}, \emph{9},
371--421.
\url{https://www.jmlr.org/papers/v9/shafer08a.html}
\bibitem{papadopoulos2002}
Papadopoulos, H.; Proedrou, K.; Vovk, V.; Gammerman, A.
Inductive Confidence Machines for Regression.
In \emph{Machine Learning: ECML 2002}; Lecture Notes in Computer Science 2430;
Springer: Berlin, 2002; pp.~345--356.
\url{https://doi.org/10.1007/3-540-36755-1_29}
\bibitem{lei2018}
Lei, J.; G'Sell, M.; Rinaldo, A.; Tibshirani, R.J.; Wasserman, L.
Distribution-Free Predictive Inference for Regression.
\emph{Journal of the American Statistical Association} \textbf{2018},
\emph{113}, 1094--1111.
\url{https://doi.org/10.1080/01621459.2017.1307116}
\bibitem{angelopoulos2023}
Angelopoulos, A.N.; Bates, S.
Conformal Prediction: A Gentle Introduction.
\emph{Foundations and Trends in Machine Learning} \textbf{2023}, \emph{16},
494--591.
\url{https://doi.org/10.1561/2200000101}
\bibitem{sadinle2019}
Sadinle, M.; Lei, J.; Wasserman, L.
Least Ambiguous Set-Valued Classifiers with Bounded Error Levels.
\emph{Journal of the American Statistical Association} \textbf{2019},
\emph{114}, 223--234.
\url{https://doi.org/10.1080/01621459.2017.1395341}
\bibitem{romano2020}
Romano, Y.; Sesia, M.; Cand\`es, E.J.
Classification with Valid and Adaptive Coverage.
In \emph{Advances in Neural Information Processing Systems 33}, 2020;
pp.~3581--3591.
\url{https://proceedings.neurips.cc/paper/2020/hash/244edd7e85dc81602b7615cd705545f5-Abstract.html}
\bibitem{tibshirani2019}
Tibshirani, R.J.; Foygel Barber, R.; Cand\`es, E.J.; Ramdas, A.
Conformal Prediction Under Covariate Shift.
In \emph{Advances in Neural Information Processing Systems 32}, 2019;
pp.~2530--2540.
\url{https://proceedings.neurips.cc/paper/2019/hash/8fb21ee7a2207526da55a679f0332de2-Abstract.html}

% --- conditional validity ---
\bibitem{vovk2012}
Vovk, V.
Conditional Validity of Inductive Conformal Predictors.
In \emph{Proceedings of the Asian Conference on Machine Learning}; PMLR 25,
2012; pp.~475--490.
\url{https://proceedings.mlr.press/v25/vovk12.html}
\bibitem{leiwasserman2014}
Lei, J.; Wasserman, L.
Distribution-Free Prediction Bands for Non-parametric Regression.
\emph{Journal of the Royal Statistical Society: Series B} \textbf{2014},
\emph{76}, 71--96.
\url{https://doi.org/10.1111/rssb.12021}
\bibitem{barber2021}
Foygel Barber, R.; Cand\`es, E.J.; Ramdas, A.; Tibshirani, R.J.
The Limits of Distribution-Free Conditional Predictive Inference.
\emph{Information and Inference} \textbf{2021}, \emph{10}, 455--482.
\url{https://doi.org/10.1093/imaiai/iaaa017}
\bibitem{vovk2003mondrian}
Vovk, V.; Lindsay, D.; Nouretdinov, I.; Gammerman, A.
Mondrian Confidence Machine.
Technical Report, Royal Holloway, University of London, 2003.
\url{https://www.alrw.net/old/04.pdf}
\bibitem{althani2026}
Althani, B.
Class-Conditional Conformal Prediction for Reliable Anomaly Detection
Under Extreme Class Imbalance.
\emph{Machine Learning and Knowledge Extraction} \textbf{2026}, \emph{8}(7), 190.
\url{https://doi.org/10.3390/make8070190}
\bibitem{geifman2017}
Geifman, Y.; El-Yaniv, R.
Selective Classification for Deep Neural Networks.
In \emph{Advances in Neural Information Processing Systems 30}, 2017;
pp.~4878--4887.
\url{https://papers.neurips.cc/paper/7073-selective-classification-for-deep-neural-networks}
\bibitem{geifman2019}
Geifman, Y.; El-Yaniv, R.
SelectiveNet: A Deep Neural Network with an Integrated Reject Option.
In \emph{Proceedings of the 36th International Conference on Machine
Learning}; PMLR 97, 2019; pp.~2151--2159.
\url{https://proceedings.mlr.press/v97/geifman19a.html}

% --- calibration ---
\bibitem{pedregosa2011}
Pedregosa, F.; Varoquaux, G.; Gramfort, A.; Michel, V.; Thirion, B.;
Grisel, O.; Blondel, M.; Prettenhofer, P.; Weiss, R.; Dubourg, V.; et al.
Scikit-learn: Machine Learning in Python.
\emph{Journal of Machine Learning Research} \textbf{2011}, \emph{12},
2825--2830.
\url{https://www.jmlr.org/papers/v12/pedregosa11a.html}
\end{thebibliography}

\end{document}

